\chapter{Introduzione}\label{chapter:introduzione} 
Negli ultimi anni la \textit{trascrittomica spaziale} (TS) ha rappresentato una delle innovazioni
più significative nel campo della biologia molecolare e della bioinformatica.  
A differenza delle tecniche di trascrittomica tradizionali, che misurano i livelli di espressione genica senza
informazioni di posizione, in quanto le cellule devono essere estratte dal tessuto, la trascrittomica spaziale consente di mappare l'attività dei geni direttamente nel contesto
del tessuto senza bisogno alcuno di doverle rimuovere.  
Questo approccio fornisce una visione bidimensionale dell'espressione genica, permettendo di studiare
non solo quali geni sono attivi, ma anche la loro posizione spaziale.

La possibilità di associare informazione spaziale all'espressione ha aperto nuove prospettive nello studio dei processi
biologici complessi, come lo sviluppo embrionale, l'organizzazione dei tessuti e la progressione di patologie come i
tumori. Tuttavia, l'analisi di questo tipo di dati presenta notevoli sfide computazionali: la quantità di dati da analizzare
è elevata e le relazioni spaziali tra geni sono intricate, e l'analisi dei cluster spaziali non sempre scova tutti i possibili pattern, spesso scambiando rumore per dati utili o viceversa.  
Da qui nasce la necessità di sviluppare strumenti e algoritmi in grado di integrare informazione geometrica, topologica di espressione al fine di trovare cluster alternativi.

\medskip
Date queste premesse si sviluppa il lavoro presentato in questa tesi, il cui obiettivo principale è stato quello di
progettare e realizzare una pipeline computazionale per l'identificazione e la classificazione delle relazioni spaziali
tra geni in un dataset 10x Visium al fine di trovare cluster spaziali alternativi.  
Partendo dai dati di espressione e dalle coordinate spaziali degli \textit{spot}, è stata costruita una lista degli spot vicini
tramite l'algoritmo dei \textit{Nearest Neighbors} (k-NN), in modo da avere le connessioni spaziali tra i vari spot del tessuto.  
Successivamente, ogni gene è stato suddiviso in \textit{sottogeni}, regioni locali di espressione contigua, mediante
l'individuazione delle \textit{componenti connesse} partendo dalle maschere binarie di espressione genica.  
Questo passaggio ha consentito di confrontare i sottogeni tra loro
e valutare diverse tipologie di rapporto spaziale: sovrapposizione, contenimento o semplice vicinanza.

\medskip
Attraverso la creazione di un algoritmo di classificazione basato su criteri quantitativi e topologici, sono state introdotte cinque categorie di relazione:
\texttt{DISGIUNTI}, \texttt{IN\_COSTA}, \texttt{A\_CONTENUTO}, \texttt{B\_CONTENUTO} e \texttt{VICINI}.  
Le informazioni ottenute sono state poi adattate in file adatti ad essere analizzati da \textit{Gephi}, software
open-source dedicato all'analisi di grafi complessi, al fine di individuare pattern spaziali di interazione tra geni che, in forma tabellare,
sarebbero stati difficilmente riconoscibili.  
In particolare, lo studio del grafo dei \texttt{VICINI} ha rivelato la presenza di cluster fortemente connessi, suggerendo
regioni di possibile co-localizzazione genica e relazioni funzionali potenzialmente rilevanti.

\medskip
Per verificare la consistenza di tali osservazioni, i risultati grafici sono stati presi e studiati con un'analisi di
\textit{clustering} condotta tramite il modello \textit{GraphST} e successivamente raffinata con \textit{Mclust}.  
L'obiettivo di questa seconda fase è valutare se le relazioni spaziali individuate tra geni corrispondono anche a pattern nei dati di espressione.  


\medskip
In sintesi, questa tesi propone un approccio completo per l'analisi della trascrittomica spaziale,
integrando strumenti di analisi dei grafi, algoritmi di clustering e visualizzazione interattiva.  
Il lavoro si inserisce nel contesto più ampio della bioinformatica computazionale e contribuisce a definire una base
metodologica utile per l'esplorazione di pattern spaziali di espressione genica, offrendo spunti concreti per futuri
sviluppi, sia sul piano tecnico che biologico.
